%% ma: L12-B14-0035
%% lop: 12
%% chuong: 5
%% bai: 14
%% dang: 12.14.19
%% loai: TN4
%% muc: TH
%% dap_an: "B"
%% nguon: "CĐ2.5.pdf giáo viên gửi 27/10/2026 (M07, Lê Minh Khôi – Nguyễn Anh Khôi, Chương II Chủ đề 5 Tích có hướng của hai vectơ), Câu 19"
%% kien_thuc: "Bài 14 (tích có hướng, phương trình mặt phẳng, khoảng cách từ điểm đến mặt phẳng), lớp 11 (thể tích khối chóp, khối hộp)"
%% trang_thai: chinh-thuc
%% ghi_chu: "Giáo viên đã duyệt 27/10/2026"
\begin{ex}
Trong không gian $Oxyz$, cho bốn điểm $A(5;1;-4)$, $B(0;1;2)$, $M(-3;-2;-4)$ và $N(4;3;5)$. Gọi $V$ là thể tích của tứ diện $ABMN$. Giá trị của $V$ nằm trong khoảng nào sau đây?
\choice{$(1;2)$}{\True $(3;4)$}{$(7;8)$}{$(5;6)$}
\loigiai{$\vec n=\left[\overrightarrow{AB},\overrightarrow{AM}\right]=(18;-48;15)$, $S_{ABM}=\dfrac12|\vec n|$, $d\left(N;(ABM)\right)=\dfrac{\left|\vec n\cdot\overrightarrow{AN}\right|}{|\vec n|}=\dfrac{21}{|\vec n|}$.

$V=\dfrac13\cdot S_{ABM}\cdot d=\dfrac{21}6=3{,}5\in(3;4)$.}
\end{ex}
